\subsection{Zielgruppe: Nachhaltigkeitsorientierte Hobbyköch:innen}\label{sec:zielgruppe}

Die Aufgabenstellung nennt Hobbyköch:innen als Zielgruppe. Die Einleitung hat bereits begründet, warum sich dieser Projektbericht auf die nachhaltigkeitsorientierte Teilgruppe verengt: Nur für Menschen, denen Resteverwertung bereits wichtig ist, ergeben Gamification-Elemente wie Rettungspunkte einen echten Anreiz.

Diese Teilgruppe lässt sich über drei Merkmale beschreiben, die weniger mit Alter oder Wohnort als mit dem Verhalten in der Küche zu tun haben. Erstens kocht sie regelmäßig selbst und plant dabei bereits mit vorhandenen Zutaten, statt ausschließlich nach Rezept einzukaufen. Zweitens ist ihr der ökologische Fußabdruck der eigenen Ernährung wichtig, auch wenn sie das nicht in jedem Gericht konsequent umsetzen kann. Drittens ist sie bereit, sich mit Gleichgesinnten über Erfolge auszutauschen, etwa in Form von Fotos gelungener Reste-Gerichte, weil dieser Austausch die eigene Motivation stützt.

Für diese Gruppe zählt neben dem Geschmack eines Rezepts auch, ob am Ende weniger im Müll landet. Daraus leiten sich zwei Anforderungen an Resteria ab: Die Plattform muss sichtbar machen, wie viel eine Person tatsächlich einspart, und einen Ort bieten, an dem dieser Erfolg von anderen wahrgenommen wird.

Ohne eigene Nutzerdaten macht eine Persona-Skizze die Zielgruppe greifbar (siehe Anhang A,~\autoref{tab:persona}). Sie fasst die drei Merkmale in einer fiktiven, aber typischen Person zusammen und dient im nächsten Abschnitt als Bezugspunkt für Konzeptentscheidungen.
